% Intended LaTeX compiler: pdflatex
\documentclass[11pt]{article}
\usepackage[utf8]{inputenc}
\usepackage[T1]{fontenc}
\usepackage{graphicx}
\usepackage{longtable}
\usepackage{wrapfig}
\usepackage{rotating}
\usepackage[normalem]{ulem}
\usepackage{amsmath}
\usepackage{amssymb}
\usepackage{capt-of}
\usepackage{hyperref}
\author{root}
\date{\today}
\title{}
\hypersetup{
 pdfauthor={root},
 pdftitle={},
 pdfkeywords={},
 pdfsubject={},
 pdfcreator={Emacs 29.3 (Org mode 9.6.15)}, 
 pdflang={English}}
\begin{document}

\tableofcontents

\section{H}
\label{sec:org8c58694}
\begin{verbatim}
1  def f():
2      return 1
\end{verbatim}

\begin{verbatim}
10  a = 1
11  b = 2
\end{verbatim}

\begin{verbatim}
16  c = 3
\end{verbatim}

\begin{verbatim}
1  (save-excursion
2    (goto-char (point-min))
\end{verbatim}
See 1 and 2 and the first.

\begin{verbatim}
1  (foo)      (a)
2  (bar)      (b)
\end{verbatim}
Ref a

\begin{verbatim}
1  (x)      (kx)
\end{verbatim}
1

\begin{verbatim}
(y)      (ly)
\end{verbatim}
ly

\begin{verbatim}
1  one
2  two
\end{verbatim}

\begin{verbatim}
 1  l1
 2  l2
 3  l3
 4  l4
 5  l5
 6  l6
 7  l7
 8  l8
 9  l9
10  l10
\end{verbatim}

\begin{verbatim}
3  (indented)      (ind)
4    (more)
\end{verbatim}
Look at line ind here.
\end{document}
